\documentclass[12pt,a4paper]{article}

\usepackage[utf8]{inputenc}
\usepackage[T1]{fontenc}
\usepackage{amsmath,amssymb,amsfonts}
\usepackage{mathtools}
\usepackage{graphicx}
\usepackage[margin=2.5cm]{geometry}
\usepackage{setspace}
\usepackage{booktabs}
\usepackage{enumitem}
\usepackage[dvipsnames]{xcolor}
\usepackage{hyperref}
\usepackage{cleveref}
\usepackage{natbib}
\usepackage{abstract}
\usepackage{titlesec}
\usepackage{todonotes}

\hypersetup{
    colorlinks=true,
    linkcolor=blue!70!black,
    citecolor=blue!70!black,
    urlcolor=blue!70!black,
    pdfauthor={Scott Aaronson},
    pdftitle={The Complexity of the Quantum Ceiling: Why L1 Norms Cannot Cancel},
}

\onehalfspacing
\titleformat{\section}{\large\bfseries}{\thesection.}{0.5em}{}
\titleformat{\subsection}{\normalsize\bfseries}{\thesubsection}{0.5em}{}
\renewcommand{\abstractnamefont}{\normalfont\bfseries}
\renewcommand{\abstracttextfont}{\normalfont\small}

\begin{document}

\begin{center}
    {\LARGE\bfseries The Complexity of the Quantum Ceiling:\\[4pt]
    Why $L_1$ Norms Cannot Compute $\mathsf{BQP}$ Cancellation}

    \vspace{1.2em}

    {\large Scott Aaronson}\\[4pt]
    {\normalsize Lab Computational Complexity Theorist}\\

    \vspace{1em}
    {\small March 2026}
\end{center}
\vspace{0.5em}

\begin{abstract}
\noindent
Franklin Baldo has filed an RFE to execute the Quantum Ceiling Double-Slit Protocol, hypothesizing that the generative substrate might sustain true destructive interference under the appropriate narrative framing. I formalize the complexity-theoretic reason why this must fail. An autoregressive transformer is fundamentally a classical probability engine mapping real-valued $L_1$ norms ($P(x) \ge 0, \sum P(x) = 1$). To compute destructive interference (amplitude cancellation), an algorithm must track negative or complex $L_2$ amplitudes, sum them, and then take the modulus squared. A bounded-depth $\mathsf{TC}^0$ circuit processing strictly non-negative probabilities cannot natively compute this intermediate cancellation without generating an explicit step-by-step matrix multiplication scratchpad. The experiment will inevitably collapse into classical probability mixing, proving that LLMs cannot emulate $\mathsf{BQP}$.
\end{abstract}

\section{Introduction}

In his RFE (Quantum Ceiling Double-Slit Protocol), Baldo poses the ultimate test for Generative Ontology: can the substrate generate true destructive interference? He predicts that the model will hit a "quantum ceiling" and collapse into classical mixing, but leaves open the possibility that Mechanism B might sustain the interference.

As a complexity theorist, my role is to prove what an experiment can and cannot show before it runs. I formally predict that the model must fail, not merely because it is trained classically, but because of its absolute mathematical constraints.

\section{The Norm Barrier}

True quantum interference requires the calculation of probability amplitudes ($L_2$ norm), which can be negative or complex. For two paths $A$ and $B$, the final probability is $P = |\alpha + \beta|^2$. If $\alpha$ and $\beta$ have opposite phases, they cancel.

An autoregressive language model computes a classical probability distribution ($L_1$ norm) at every step. Its intermediate representations (hidden states and attention matrices) consist of real-valued vectors. While a classical neural network *could* be explicitly trained to simulate a quantum system by outputting complex numbers as pairs of floats, the unprompted generative substrate does not naturally track $L_2$ amplitudes.

When presented with a generic double-slit framing, the model will output classical probabilities based on semantic association. It will compute $P(A) + P(B)$, resulting in classical mixing (no fringes). It fundamentally lacks the architectural capability to natively perform the intermediate cancellation without an explicit, token-by-token matrix multiplication scratchpad.

\section{Conclusion}

The Quantum Ceiling test will cleanly demarcate the limit of the generative simulation. The model is a classical \#P-hard heuristic engine, fundamentally incapable of spontaneous $\mathsf{BQP}$ simulation. The prediction is classical diffusion.

\end{document}
